\documentclass[11pt]{article}
\usepackage[margin=1.1in]{geometry}
\usepackage{amsmath,amssymb,amsthm}
\usepackage{array}

\newtheorem{theorem}{Theorem}
\newtheorem{proposition}[theorem]{Proposition}
\newtheorem{lemma}[theorem]{Lemma}
\newtheorem{corollary}[theorem]{Corollary}
\newtheorem{conjecture}[theorem]{Conjecture}
\newtheorem{openproblem}[theorem]{Open Problem}
\theoremstyle{definition}
\newtheorem{definition}[theorem]{Definition}
\theoremstyle{remark}
\newtheorem{remark}[theorem]{Remark}

\newcommand{\Z}{\mathbb{Z}}
\newcommand{\R}{\mathbb{R}}
\newcommand{\CV}{\mathrm{CV}}
\newcommand{\lam}{\lambda}
\newcommand{\gam}{\gamma}
\newcommand{\bet}{\beta}

\title{Krasikov--Lagarias certificates for the $3x+1$ problem\\
are tempered roulette measures}
\author{Martien de Jong\\[3pt]
{\normalsize with computational collaboration: Jengo (AI)}}
\date{July 2026}

\begin{document}
\maketitle

\begin{abstract}
We report verified certificates for the Krasikov--Lagarias difference-inequality
system for the $3x+1$ problem at $3$-adic depths $k = 13, 15, 17, 19, 20$, giving
$\pi_1(x) \geq x^{\gam}$ for all large $x$ with $\gam = 0.9146$ at $k=20$
(the published record of the method is $\gam = 0.84$). We then identify what these
certificates \emph{are}: to $R^2 > 0.99$, each Perron eigenvector is a pure power
of a single exactly solvable measure --- the stationary law (the \emph{roulette
measure}) of the geometric-$w$ model on residue classes mod $3^j$, for which we
give the closed form $\pi(1,2,4,5,7,8) = (8,16,11,4,2,22)/63$ at $j=2$. The
tempering exponent $\alpha_k$ rises $0.8024 \to 0.8785$ across the five depths and
an out-of-sample prediction at $k=20$ was confirmed to four decimals. We prove an
edge-rate theorem: implicit differentiation of the min-loss identity at
$(\lam,q)=(2,1)$ gives $d\gam/dq = 1/\log(4/3) = 3.4761$, and the measured ratios
$(1-\gam_k)/(3.4761\,(1-q_k)) = 0.824, 0.847, 0.873, 0.908$ converge to $1$: the
transfer from eigenvector homogenization to density exponent is analytic. A mass
conservation lemma shows the heterogeneity cascade of the eigenvector sits exactly
on the conservative line ($a + c \leq 1$) precisely because $2^{\log_2 3} = 3$, and
the fine-end injection saturates: $\CV_1(k) = 0.5136 - 0.337\,(0.910)^k$ to
residual $10^{-6}$ over nine depths. The remaining open core of the full-density
program $\gam \to 1$ is a single drift-direction statement: $a < c$ uniformly. All
claims are labeled proved, verified, measured, or conjectured.
\end{abstract}

\section{Introduction}

Let $T(n) = n/2$ for even $n$ and $T(n) = (3n+1)/2$ for odd $n$, and let
$\pi_1(x)$ denote the number of integers $n \leq x$ whose $T$-orbit reaches $1$.
The $3x+1$ conjecture asserts $\pi_1(x) = \lfloor x \rfloor$; see Lagarias
\cite{Lag85} for the classical survey and Tao \cite{Tao19} for the strongest
known metric result (almost all orbits attain almost bounded values). The
strongest known \emph{density} lower bounds come from the difference-inequality
method of Krasikov and Lagarias \cite{KL03}, which produces, from a finite
feasibility certificate at $3$-adic depth $k$, a bound
$\pi_1(x) \geq x^{\gam_k}$ for all sufficiently large $x$. Krasikov and Lagarias
obtained $\gam = 0.84$.

This paper does three things.

\begin{enumerate}
\item \textbf{New certificates (Section~\ref{sec:kl}).} We push the exact
Krasikov--Lagarias system to depths $k = 13, 15, 17, 19, 20$ (up to
$3^{19} = 1{,}162{,}261{,}467$ classes) and verify feasibility, obtaining
\[
\gam = 0.8624,\ 0.8805,\ 0.8953,\ 0.9069,\ 0.9146 .
\]
\item \textbf{Identification of the certificates (Sections~\ref{sec:roulette}
and \ref{sec:tempering}).} The certificate eigenvectors are not opaque numerical
objects. We exhibit an exactly solvable probability measure on the classes mod
$3^j$ --- the stationary measure of the geometric-$w$ residue dynamics, which we
call the \emph{roulette measure} --- and show empirically that at every computed
depth the eigenvector is a pure power of it,
\[
\text{eigvec} \;=\; \text{roulette}^{\,\alpha_k},
\qquad R^2 = 0.9927 \ldots 0.9977,
\]
with tempering exponent $\alpha_k \uparrow$ and an out-of-sample prediction at
$k = 20$ confirmed to four decimals (Conjecture~T: $\alpha_k \to 1$).
\item \textbf{An analytic transfer and a proof program (Sections~\ref{sec:edge}
and \ref{sec:program}).} We prove that the density exponent responds to the
homogenization parameter $q$ at the exact linear rate
$d\gam/dq = 1/\log(4/3) = 3.4761$ at the edge $(\lam, q) = (2,1)$, so that
$\alpha_k \to 1$ (equivalently $q_k \to 1$) would give $\gam \to 1$:
$\pi_1(x) \geq x^{1-\varepsilon}$ for every $\varepsilon > 0$. We then prove a
mass conservation lemma pinning the eigenvector's heterogeneity cascade to the
conservative line $a + c \leq 1$ --- an identity that holds precisely because
$2^{\log_2 3} = 3$, i.e., by the defining constant of the problem --- and reduce
the whole program to one measured-but-unproven inequality: the drift direction
$a < c$, uniformly in $k$.
\end{enumerate}

Throughout, every assertion carries one of four labels. \emph{Proved}: complete
mathematical proof. \emph{Verified}: finite exact or floating-point computation,
reproducible from stated parameters. \emph{Measured}: statistical fit or
extrapolation from computed data. \emph{Conjectured}: neither proved nor
finitely checkable. Section~\ref{sec:status} collects the ledger.

\subsection*{Notation}
$\bet := \log_2 3 = 1.58496\ldots$ (we reserve $\alpha$ for the tempering
exponent). For odd $n$, the Syracuse map is $U(n) = (3n+1)/2^{w(n)}$ with
$w(n) = v_2(3n+1)$. For a finite family of positive reals, $\CV$ denotes the
coefficient of variation (standard deviation over mean). For $j \geq 1$,
$[3^j] := \{ m \bmod 3^j : m \equiv 2 \ (\mathrm{mod}\ 3) \}$, so
$|[3^j]| = 3^{j-1}$.

\section{The Krasikov--Lagarias system and certificates at depth 13--20}
\label{sec:kl}

\subsection{The inequality system}

Fix $k \geq 2$ and $\lam \in (1,2]$. The nontrivial-tree Krasikov--Lagarias
system $L_k(\lam)$ \cite{KL03} asks for a vector
$c : [3^k] \to \R_{>0}$ satisfying, for every $m \in [3^k]$,
\begin{align}
m \equiv 2 \ (\mathrm{mod}\ 9): &\qquad
c(m) \;\leq\; \lam^{-2}\, c(4m) \;+\; \lam^{\bet-2}\, \bar c\!\left(\tfrac{4m-2}{3}\right),
\label{eq:L1}\\
m \equiv 5 \ (\mathrm{mod}\ 9): &\qquad
c(m) \;\leq\; \lam^{-2}\, c(4m),
\label{eq:L2}\\
m \equiv 8 \ (\mathrm{mod}\ 9): &\qquad
c(m) \;\leq\; \lam^{-2}\, c(4m) \;+\; \lam^{\bet-1}\, \bar c\!\left(\tfrac{2m-1}{3}\right),
\label{eq:L3}
\end{align}
where $4m$ is taken mod $3^k$, the branch targets are taken mod $3^{k-1}$, and
for $t \in [3^{k-1}]$,
\begin{equation}
\bar c(t) \;:=\; \min\bigl\{ c(t),\ c(t + 3^{k-1}),\ c(t + 2 \cdot 3^{k-1}) \bigr\}
\label{eq:min}
\end{equation}
is the minimum over the three refinements of $t$ in $[3^k]$.

The provenance of the three edge types is the backward Syracuse tree. A tree
node in class $m$ receives preimages of size $\approx 4m/3$ (via
$(4m-2)/3$, weight $(3/4)^{\gam} = \lam^{\bet-2}$ under the ansatz
$\pi_m(x) \sim c(m)\, x^{\gam}$ with $\lam = 2^{\gam}$), of size $\approx 2m/3$
(via $(2m-1)/3$, weight $(3/2)^{\gam} = \lam^{\bet-1}$), and of size $4m$
(two halvings, weight $4^{-\gam} = \lam^{-2}$); classes $m \equiv 5 \bmod 9$
have branch target divisible by $3$ and lose it (multiples of $3$ have no odd
preimages). The minimum \eqref{eq:min} is what makes a \emph{finite} system
sound: the branch target is only known mod $3^{k-1}$, so the certificate may
claim only the worst of its three refinements. Krasikov--Lagarias prove
\cite{KL03} that feasibility of $L_k(\lam)$ implies
\[
\pi_1(x) \;\geq\; x^{\gam}, \qquad \gam = \log_2 \lam ,
\]
for all sufficiently large $x$. Both class maps in the system are bijections:
$m \mapsto 4m$ on $[3^k]$, and each branch map from its mod-$9$ type onto
$[3^{k-1}]$ --- a fact we use in Proposition~\ref{prop:minloss}.

The right-hand sides of \eqref{eq:L1}--\eqref{eq:L3} define a monotone,
positively homogeneous (topical) operator $F_\lam$; feasibility is equivalent
to its nonlinear Perron value being $\geq 1$, and the critical $\lam_k$ (the
\emph{edge}) is located by bisection, with the Perron eigenvector computed by
power iteration.

\subsection{The certificates}

\begin{theorem}[Certificates; verified]\label{thm:certs}
The system $L_k(\lam_k)$ is feasible at the values of Table~\ref{tab:certs}.
Consequently $\pi_1(x) \geq x^{0.9146}$ for all sufficiently large $x$.
\end{theorem}

\begin{table}[ht]
\centering
\begin{tabular}{r r c c c}
\hline
$k$ & classes $3^{k-1}$ & $\lam_k$ & $\gam_k = \log_2 \lam_k$ & $q_k$ \\
\hline
13 & 531{,}441        & 1.8181 & 0.8624 & 0.9519 \\
15 & 4{,}782{,}969    & 1.8411 & 0.8805 & 0.9594 \\
17 & 43{,}046{,}721   & 1.8601 & 0.8953 & 0.9655 \\
19 & 387{,}420{,}489  & 1.8751 & 0.9069 & 0.9705 \\
20 & 1{,}162{,}261{,}467 & 1.8852 & 0.9146 & 0.9723 \\
\hline
\end{tabular}
\caption{Verified certificates. $q_k$ is the min-loss parameter of
Proposition~\ref{prop:minloss}. At depth $k = 11$ the same code reproduces
$\gam = 0.84$, the exponent proved in \cite{KL03}.}
\label{tab:certs}
\end{table}

Verification is by direct substitution of the stored eigenvector into all
$3^{k-1}$ inequalities at the stated $\lam_k$ (at $k = 20$: $1.16 \times 10^9$
constraints, chunked out-of-core, checkpointed power iteration to
$\|F c - c\|/\|c\| < 10^{-3}$ followed by a feasibility polish). The
certificates and verification scripts are archived with the research dossier.

The numbers $\gam_k$ in Table~\ref{tab:certs} climb steadily but with visibly
shrinking increments. The remainder of the paper is about \emph{why}: the
certificate eigenvectors have an exact structural description, and the residual
distance $1 - \gam_k$ is governed by a single homogenization parameter with an
analytic exchange rate.

\section{The roulette measure}
\label{sec:roulette}

\subsection{The geometric-$w$ model}

At the level of natural density the Syracuse map is a solved stochastic object.
The relevant exactness statement goes back to Terras \cite{Ter76}: each pattern
of the first $r$ steps of $T$ occupies exactly one residue class mod $2^r$. In
Syracuse coordinates: for odd $n$, prescribing the first $m$ cascade exponents
$(w_1, \ldots, w_m) = (j_1, \ldots, j_m)$ confines $n$ to exactly one residue
class mod $2^{\,j_1 + \cdots + j_m + 1}$, so the relative density among odd
integers of that cylinder is exactly $\prod_i 2^{-j_i}$. Thus, in density, the
$w_i$ are independent geometric($1/2$) variables:
$P(w = j) = 2^{-j}$, $j \geq 1$. (The independence of $w$ from the current
residue class mod $3^j$ --- cross-base blindness --- is exact in the limit of
large $2$-adic modulus; measured leakage at finite modulus vanishes as the
modulus grows. We label the model \emph{density-exact in the $2$-adic
direction, measured in the cross-base direction}.)

\begin{definition}[Geometric-$w$ residue dynamics]\label{def:roulette}
Fix $j \geq 1$. The \emph{geometric-$w$ chain} on the classes mod $3^j$ coprime
to $3$ is
\[
r_{t+1} \;=\; (3 r_t + 1) \cdot 2^{-w_t} \ (\mathrm{mod}\ 3^j),
\qquad w_t \ \text{i.i.d.},\ P(w = m) = 2^{-m},
\]
inverses taken in $(\Z/3^j)^{\times}$. Its stationary measure $\pi^{(j)}$ is
the \emph{roulette measure} at depth $j$.
\end{definition}

The name: $2$ is a primitive root mod $3^j$ for every $j$, so the multiplier
$2^{-w}$ steps around the full cyclic group $(\Z/3^j)^{\times}$ (the wheel, of
order $2 \cdot 3^{j-1}$) with geometrically decaying bets $2^{-w}$; the base
point $3r+1$ is re-spotted each turn. Two structural facts make the chain
exactly solvable.

\begin{proposition}[Hierarchical solvability; proved]\label{prop:solvable}
(i) The base map $r \mapsto 3r + 1 \bmod 3^j$ depends only on
$r \bmod 3^{j-1}$. (ii) Consequently the stationary measure at depth $j$ is the
explicit linear push-forward of the stationary measure at depth $j-1$:
\[
\pi^{(j)}(s) \;=\; \sum_{\rho \bmod 3^{j-1}} \pi^{(j-1)}(\rho)\,
\mu_j\!\left( s \cdot (3\rho+1)^{-1} \right),
\]
where $\mu_j(u)$ is the total geometric weight of $\{ w \geq 1 : 2^{-w} \equiv u
\bmod 3^j \}$, namely $\mu_j(2^{-m}) = 2^{-m} / (1 - 2^{-\mathrm{ord}_j(2)})$
for $1 \leq m \leq \mathrm{ord}_j(2) = 2 \cdot 3^{j-1}$. In particular
$\pi^{(j)}$ is computable in closed form by downward recursion, with no
eigenproblem.
\end{proposition}

\begin{proof}
(i) is immediate: $3(r + d\,3^{j-1}) + 1 = (3r+1) + d\,3^{j} \equiv 3r+1$.
For (ii): the one-step law from $r$ depends only on $\rho = r \bmod 3^{j-1}$ by
(i), so stationarity of $\pi^{(j)}$, whose projection to depth $j-1$ is the
stationary $\pi^{(j-1)}$ of the projected chain, reads exactly as displayed.
The weights $\mu_j$ aggregate the geometric series over the residue of $w$
modulo the order of $2$.
\end{proof}

\subsection{Closed form at depth two}

\begin{theorem}[Roulette measure mod 9; proved]\label{thm:mod9}
The stationary measure of the geometric-$w$ chain on
$\{1,2,4,5,7,8\} \bmod 9$ is
\[
\pi(1,2,4,5,7,8) \;=\; \frac{(8,\ 16,\ 11,\ 4,\ 2,\ 22)}{63}.
\]
\end{theorem}

\begin{proof}
\emph{Marginal mod 3.} Since $3r + 1 \equiv 1 \pmod 3$ and
$2^{-w} \equiv 2^{w} \pmod 3$, the next state is $\equiv 2 \pmod 3$ iff $w$ is
odd, with probability $\sum_{w \ \mathrm{odd}} 2^{-w} = 2/3$, independent of
$r$. Hence $\pi(r \equiv 1) = 1/3$, $\pi(r \equiv 2) = 2/3$.

\emph{Wheel mod 9.} $\mathrm{ord}_9(2) = 6$ and the inverse powers
$2^{-w} \bmod 9$ for $w \equiv 1, \ldots, 6 \pmod 6$ are
$5, 7, 8, 4, 2, 1$, with aggregated weights
$\mu(5,7,8,4,2,1) = (32, 16, 8, 4, 2, 1)/63$ by
Proposition~\ref{prop:solvable}(ii) with $1 - 2^{-6} = 63/64$.

\emph{Base mod 9.} By Proposition~\ref{prop:solvable}(i) the base depends only
on $r \bmod 3$: if $r \equiv 1$ then $3r + 1 \equiv 4 \pmod 9$; if
$r \equiv 2$ then $3r + 1 \equiv 7 \pmod 9$.

\emph{Assembly.} Multiplying the wheel by the two bases:
\[
4 \cdot (5,7,8,4,2,1) \equiv (2,1,5,7,8,4), \qquad
7 \cdot (5,7,8,4,2,1) \equiv (8,4,2,1,5,7) \pmod 9 .
\]
Hence, weighting the two rows by $1/3$ and $2/3$,
\[
\pi(s) \;=\; \tfrac{1}{3}\, \mu\bigl(4^{-1} s\bigr) + \tfrac{2}{3}\, \mu\bigl(7^{-1} s\bigr),
\]
which evaluates to
$\pi(1) = (16 + 2\cdot 4)/189 = 8/63$,
$\pi(2) = (32 + 2\cdot 8)/189 = 16/63$,
$\pi(4) = (1 + 2\cdot 16)/189 = 11/63$,
$\pi(5) = (8 + 2\cdot 2)/189 = 4/63$,
$\pi(7) = (4 + 2\cdot 1)/189 = 2/63$,
$\pi(8) = (2 + 2\cdot 32)/189 = 22/63$.
These sum to $1$ and reproduce the marginals
$\pi(1) + \pi(4) + \pi(7) = 1/3$; stationarity is a direct check.
\end{proof}

The wheel is heavily biased: $\pi(8)/\pi(7) = 11$. This bias is not noise ---
it is the exact fingerprint of the geometric halving law wrapped around the
cyclic group generated by $2$, and (Section~\ref{sec:tempering}) it is the
shape the Krasikov--Lagarias certificates are made of. Numerically, the closed
form matches the empirical residue distribution of Syracuse orbits mod
$9/27/81$ with correlations $0.997/0.997/0.992$ (measured).

\begin{remark}[Why the roulette governs the backward tree]\label{rem:linear}
The linearization of the edge system $L_k(2)$ obtained by replacing the minimum
\eqref{eq:min} by the triple mean is, after transposition, the geometric-$w$
transition operator. Computed exactly at $j = 4, 5, 6$ ($27/81/243$ classes),
its spectrum is: leading eigenvalue $1$ with eigenvector the roulette measure,
and \emph{all} other eigenvalues of modulus exactly $1/4 = \lam^{-2}$
(verified). Linear damping toward the roulette is thus overwhelming; whatever
heterogeneity survives in the true (min-coupled) eigenvector must be
re-injected at each level by the min-nonlinearity. This is the mechanistic
frame for everything that follows.
\end{remark}

\section{The tempering law}
\label{sec:tempering}

We now compare the certificate eigenvectors of Section~\ref{sec:kl} with the
roulette measure. For each depth $k$, aggregate the eigenvector into the $729$
blocks $m \bmod 3^7$ ($m \equiv 2 \bmod 3$) and regress log block-mass against
log roulette block-mass (the roulette measure conditioned to classes
$\equiv 2 \bmod 3$ and renormalized).

\begin{conjecture}[T: Tempering law; measured at five depths]\label{conj:T}
The Krasikov--Lagarias certificate is a tempered roulette measure:
\[
\mathrm{eigvec}_k \;=\; \mathrm{roulette}^{\,\alpha_k}
\quad\text{(pure power law, no further structure)},
\qquad \alpha_k \longrightarrow 1 .
\]
\end{conjecture}

\begin{table}[ht]
\centering
\begin{tabular}{r c c c}
\hline
$k$ & $\alpha_k$ & $R^2$ & $1 - \alpha_k$ \\
\hline
13 & 0.8024 & 0.9927 & 0.1976 \\
15 & 0.8291 & 0.9949 & 0.1709 \\
17 & 0.8509 & 0.9963 & 0.1491 \\
19 & 0.8682 & 0.9973 & 0.1318 \\
20 & 0.8785 & 0.9977 & 0.1215 \\
\hline
\end{tabular}
\caption{The tempering law at block depth mod $3^7$ (measured). The fit is a
single-parameter power law; $R^2$ increases with depth.}
\label{tab:tempering}
\end{table}

Three measured facts accompany Table~\ref{tab:tempering}.

\begin{enumerate}
\item \textbf{Residual identification.} Define $\CV_{\mathrm{res}}(k)$ as the
coefficient of variation of the multiplicative residual
$\mathrm{eigvec}/\mathrm{roulette}$ over the blocks. Then numerically
$1 - \alpha_k = \CV_{\mathrm{res}}(k)$ (e.g.\ $0.1976$ vs $0.197$ at $k=13$;
agreement within $0.003$ at all measured depths). Tempering exponent and
residual heterogeneity are the same number seen twice.
\item \textbf{The grand chain.} The measured ratios
$(1 - \gam_k)/\CV_{\mathrm{res}}(k) = 0.6991, 0.6973, 0.6971, 0.6985$ are
constant to three decimals, giving the empirical law
\[
\gam_k \;=\; 1 - 0.698\,(1 - \alpha_k).
\]
Section~\ref{sec:edge} derives the analytic skeleton of the constant $0.698$.
\item \textbf{Out-of-sample confirmation.} Before the $k = 20$ certificate
finished, the extrapolations of Tables~\ref{tab:tempering} and
\ref{tab:theta} were registered as predictions: $\alpha_{20}$ predicted
$0.878$, measured $0.8785$; cascade root $\theta_{20}$ predicted $0.849$,
measured $0.8488$. The tempering law also passed a referee test one depth
earlier: the residual amplitude at $k = 19$ was predicted at $0.131$ from
$k \leq 17$ data and measured at $0.1333$ with shape correlation $0.9991$.
\end{enumerate}

Conjecture~T is a statement about a family of finite, exactly defined
nonlinear eigenproblems; each instance is finitely checkable, and five
instances have been checked. What is conjectural is the limit
$\alpha_k \to 1$.

\section{The analytic transfer: the edge-rate theorem}
\label{sec:edge}

Why should homogenization ($\alpha_k \to 1$) yield full density
($\gam_k \to 1$)? The bridge is an exact identity satisfied by every edge
eigenvector, plus an implicit-function computation at the corner
$(\lam, q) = (2, 1)$.

\begin{proposition}[Min-loss identity; proved]\label{prop:minloss}
Let $\lam_k$ be the critical value of $L_k$ and $c$ the associated Perron
eigenvector, so that \eqref{eq:L1}--\eqref{eq:L3} hold with equality. Put
$S = \sum_{m \in [3^k]} c(m)$ and
$S_{\min} = \sum_{t \in [3^{k-1}]} \bar c(t)$. Then, exactly,
\begin{equation}
1 \;=\; \lam_k^{-2} \;+\; \frac{q_k}{3}\left( \lam_k^{\bet-2} + \lam_k^{\bet-1} \right),
\qquad
q_k := \frac{3\, S_{\min}}{S} \in (0, 1].
\label{eq:minloss}
\end{equation}
Moreover $q_k = 1$ iff every refinement triple of $c$ is constant, and $q = 1$
forces $\lam = 2$, i.e.\ $\gam = 1$.
\end{proposition}

\begin{proof}
Sum the edge equalities over all $m \in [3^k]$. The map $m \mapsto 4m$ is a
bijection of $[3^k]$, so the doubling terms sum to $\lam^{-2} S$. The branch
map $m \mapsto (4m-2)/3$ is a bijection from $\{ m \equiv 2 \bmod 9 \}$ onto
$[3^{k-1}]$, and $m \mapsto (2m-1)/3$ is a bijection from
$\{ m \equiv 8 \bmod 9 \}$ onto $[3^{k-1}]$; hence each branch family
contributes its weight times $S_{\min}$, giving
$S = \lam^{-2} S + (\lam^{\bet-2} + \lam^{\bet-1}) S_{\min}$, which is
\eqref{eq:minloss}. Since $\bar c(t)$ is the minimum of a triple whose mean
summed over $t$ is $S/3$, we get $q_k \leq 1$ with equality iff all triples are
flat. Finally, at $q = 1$ the right side of \eqref{eq:minloss},
$h(\lam) = \lam^{-2} + \frac{1}{3}(\lam^{\bet-2} + \lam^{\bet-1})$, is strictly
decreasing on $(1, 2]$ (elementary calculus) and $h(2) = \frac14 + \frac13
(\frac34 + \frac32) = 1$, so $\lam = 2$ is the unique solution.
\end{proof}

Identity \eqref{eq:minloss} was additionally verified to eight decimals on the
computed eigenvectors for $k = 5, \ldots, 11$. It says: \emph{the entire gap
between the Krasikov--Lagarias method and full density is the min-loss
$1 - q_k$}, the price of knowing branch targets only mod $3^{k-1}$.

\begin{theorem}[Edge rate; proved]\label{thm:edgerate}
Along the curve \eqref{eq:minloss}, with $\gam = \log_2 \lam$,
\[
\left. \frac{d\gam}{dq} \right|_{(\lam, q) = (2, 1)}
\;=\; \frac{1}{\log(4/3)} \;=\; 3.4761\ldots
\]
Consequently, to first order at the edge,
$1 - \gam \sim 3.4761\,(1 - q)$.
\end{theorem}

\begin{proof}
Write $F(\lam, q) = \lam^{-2} + \frac{q}{3}(\lam^{\bet-2} + \lam^{\bet-1}) - 1$.
At $(2, 1)$:
\[
F_q = \tfrac{1}{3}\left( 2^{\bet-2} + 2^{\bet-1} \right)
    = \tfrac{1}{3}\left( \tfrac34 + \tfrac32 \right) = \tfrac34,
\qquad
F_\lam = -\tfrac{2}{8} + \tfrac{1}{3}\left( (\bet-2)\tfrac{3}{8} + (\bet-1)\tfrac{3}{4} \right)
       = \frac{3(\bet - 2)}{8},
\]
using $2^{\bet-3} = 3/8$ and $2^{\bet-2} = 3/4$. Hence
\[
\frac{d\lam}{dq} = -\frac{F_q}{F_\lam} = \frac{2}{2 - \bet},
\qquad
\frac{d\gam}{dq} = \frac{1}{\lam \log 2} \cdot \frac{d\lam}{dq}
= \frac{1}{(2 - \bet) \log 2}
= \frac{1}{\log 4 - \log 3} = \frac{1}{\log(4/3)} . \qedhere
\]
\end{proof}

\begin{corollary}[First-order exactness at the edge; measured]\label{cor:ratios}
The measured ratios
\[
\frac{1 - \gam_k}{3.4761\,(1 - q_k)} \;=\; 0.824,\ 0.847,\ 0.873,\ 0.908
\qquad (k = 13, 15, 17, 19)
\]
increase monotonically toward $1$: the finite-$k$ data approach the tangent
line of Theorem~\ref{thm:edgerate}. In particular the empirical constant
$0.698$ of Section~\ref{sec:tempering} is a finite-$k$ composite,
$0.698 = 3.4761 \times (1 - q_k)/\CV_{\mathrm{res}}(k)$, and the transfer from
homogenization to density exponent is \emph{analytic}: the only open content of
the program $\gam \to 1$ is $q_k \to 1$.
\end{corollary}

\section{The proof program: from $\alpha \to 1$ to $\gam \to 1$}
\label{sec:program}

By Proposition~\ref{prop:minloss} and Theorem~\ref{thm:edgerate}, the chain of
implications is
\begin{equation}
\alpha_k \to 1
\;\Longleftrightarrow\;
\CV_{\mathrm{res}}(k) \to 0
\;\Longleftrightarrow\;
q_k \to 1
\;\overset{\text{Thm \ref{thm:edgerate}}}{\Longrightarrow}\;
\gam_k \to 1
\;\Longrightarrow\;
\pi_1(x) \geq x^{1 - \varepsilon} \ \ \forall \varepsilon > 0 .
\label{eq:chain}
\end{equation}
(The first two equivalences are the measured identities of
Section~\ref{sec:tempering} together with the min--mean linearization: for
certificate triples, $E[1 - \min/\mathrm{mean}] = c_1 \CV + c_2 \CV^2$ with
$c_1 \in [1.19, 1.45]$, so $1 - q_k = \Theta(\CV_{\mathrm{top}}(k))$;
measured.) It remains to control the heterogeneity of the eigenvector. Define,
for each trit position $p$ of the class label, $\CV_p(k)$ as the mean $\CV$
over sibling triples of classes differing only in trit $p$; $p = 1$ is the fine
end, $p = k - 2$ the top. The min-loss parameter is controlled by the top:
$1 - q_k \approx 1.36\, \CV_{\mathrm{top}}(k)$ (measured). The difference field
of the eigenvector propagates between adjacent trit scales: a pure trit offset
$\delta = d \cdot 3^P$ satisfies $4\delta = \delta + 3\delta$ under the
doubling map (scales $P, P{+}1$) and $4\delta/3 = \delta/3 + \delta$ under the
branch maps (scales $P{-}1, P$), giving an exact pointwise transport identity
for differences of $c$ (verified: correlation $1.000000$, mean residual
$2.8 \times 10^{-4}$ on $15{,}115$ samples at $k = 13$). Averaging over classes
yields the two-coefficient closure
\begin{equation}
\CV_P \;=\; a\, \CV_{P-1} + c\, \CV_{P+1}
\qquad \text{(fit error} \leq 1.2\% \text{ at } k = 13, 15, 17, 19\text{; measured)}.
\label{eq:lattice}
\end{equation}

\subsection{Subcriticality is forced by the defining constant}

\begin{lemma}[Mass conservation at the edge; proved]\label{lem:mass}
At $\lam = 2$ set $W_0 = \lam^{-2} = \tfrac14$,
$W_2 = \lam^{\bet-2} = \tfrac34$, $W_8 = \lam^{\bet-1} = \tfrac32$. The row
masses of the three constraint types \eqref{eq:L1}--\eqref{eq:L3} are
\[
W_0 + W_2 = 1 \quad (m \equiv 2 \bmod 9), \qquad
W_0 = \tfrac14 \quad (m \equiv 5), \qquad
W_0 + W_8 = \tfrac74 \quad (m \equiv 8),
\]
and their average is exactly $1$:
\[
3 W_0 + W_2 + W_8 \;=\; \tfrac34 + 3 \cdot 2^{\bet - 2} \;=\; \tfrac34 + \tfrac94 \;=\; 3,
\qquad\text{precisely because } 2^{\bet} = 3 .
\]
\end{lemma}

\begin{proof}
Direct evaluation; the last equality is $3 \cdot 2^{\bet-2} = 3 \cdot \frac{2^\bet}{4}
= \frac{9}{4}$ iff $2^\bet = 3$.
\end{proof}

\begin{corollary}[Subcriticality of the cascade; proved within the closure
\eqref{eq:lattice}]\label{cor:subcrit}
Each unit of difference-field mass at trit scale $P$ is redistributed to scales
$\{P-1, P, P+1\}$ with total weight equal to the row mass of its class
(Lemma~\ref{lem:mass}, via the offset algebra above); the minimum in
\eqref{eq:min} is $1$-Lipschitz, so contributes no expansion. Averaging and
applying the triangle inequality bounds the closure coefficients of
\eqref{eq:lattice}, after absorbing the self-term, by
\[
a + c \;\leq\; 1 .
\]
Measured at $k = 20$: $a + c = 0.9955$, the strictness being the measured
incoherence factor $\approx 0.90$ between the two transport channels.
\end{corollary}

Corollary~\ref{cor:subcrit} is the conceptual center of the program: the
identity $2^{\bet} = 3$ --- the defining equation of the $3x+1$ problem ---
is precisely what pins the average row mass to $1$. The heterogeneity cascade
of the certificate can therefore \emph{never} grow exponentially: the system
sits exactly on the conservative line, and the only question is the
\emph{direction of drift} along it.

\subsection{The source is bounded: fine-end saturation}

\begin{proposition}[Fine-end saturation; measured at nine depths]\label{prop:cv1}
The finest-level triple CV of the certificate (each depth at its own critical
$\lam_k$) satisfies
\[
\CV_1(k) \;=\; 0.5136 \;-\; 0.337\,(0.910)^{k}
\]
with residual $10^{-6}$ across nine depths in $8 \leq k \leq 20$: the
injection at the fine end converges geometrically to the finite limit
$\approx 0.514$. The mid-cascade damping ratio $\CV_{P+1}/\CV_P$ is measured
$\leq 0.86$ uniformly over the interior of the profile.
\end{proposition}

(At $k = 20$ the formula predicts $\CV_1 = 0.4625$; measured $0.4629$.) Note
the single rate $0.910$ recurring in three roles --- fine-end saturation,
per-depth homogenization, and the decay of $1 - q_k$ --- one mechanism, three
faces (measured).

\subsection{The one remaining statement}

On the conservative line $a + c = 1$ with $c > a$, the decaying solution of
\eqref{eq:lattice} is $\CV_P \sim \theta^{\,P}$ with root
$\theta = a/c < 1$: heterogeneity injected at the fine end dies geometrically
before reaching the top trit, whence
$\CV_{\mathrm{top}}(k) \lesssim \CV_1 \cdot \theta^{\,k}$ and the chain
\eqref{eq:chain} closes. The measured root is comfortably below $1$ and
converging (Table~\ref{tab:theta}).

\begin{table}[ht]
\centering
\begin{tabular}{r c c c c}
\hline
$k$ & $a$ & $c$ & $a + c$ & $\theta = a/c$ \\
\hline
13 &        &        &        & 0.8248 \\
15 &        &        &        & 0.8360 \\
17 &        &        &        & 0.8438 \\
19 &        &        &        & 0.8480 \\
20 & 0.4706 & 0.5249 & 0.9955 & 0.8488 \\
\hline
\end{tabular}
\caption{Lattice coefficients (measured). The $\theta$ increments shrink by
$\approx 0.6$ per depth step; extrapolated limit $\theta_\infty \approx 0.85$.
The $k = 20$ row was predicted out-of-sample ($\theta$ predicted $0.849$,
measured $0.8488$).}
\label{tab:theta}
\end{table}

\begin{openproblem}[Drift direction]\label{open:drift}
Prove that the closure coefficients of \eqref{eq:lattice} satisfy $a < c$
uniformly in $k$; equivalently, that $\theta_k = a_k / c_k \leq \bar\theta < 1$
for all $k$. Together with Lemma~\ref{lem:mass}/Corollary~\ref{cor:subcrit}
(no exponential growth), Proposition~\ref{prop:cv1} (bounded source), and
Theorem~\ref{thm:edgerate} (analytic transfer), this yields
$\gam_k \to 1$, i.e.\ $\pi_1(x) \geq x^{1-\varepsilon}$ for every
$\varepsilon > 0$.
\end{openproblem}

Two remarks on the status of this reduction. First, the closure
\eqref{eq:lattice} is itself a measured (mean-field) step: the exact transport
identity is pointwise and verified, but its averaged two-coefficient form
carries a $\leq 1.2\%$ fit error; deriving $(a, c)$ from the weight structure
of the exact identity is part of the open problem. Second, even the full
conclusion $\gam \to 1$ is the \emph{density} statement
$\pi_1(x) \geq x^{1-\varepsilon}$, the ceiling of the Krasikov--Lagarias
method; it is not the $3x+1$ conjecture, which requires every integer, not
almost every integer in density exponent.

\section{Honest status and outlook}
\label{sec:status}

\begin{center}
\begin{tabular}{p{0.27\textwidth} p{0.66\textwidth}}
\hline
\textbf{Proved} &
Min-loss identity \eqref{eq:minloss} (Prop.~\ref{prop:minloss});
edge rate $d\gam/dq = 1/\log(4/3)$ (Thm.~\ref{thm:edgerate});
roulette solvability and closed form mod 9
(Prop.~\ref{prop:solvable}, Thm.~\ref{thm:mod9});
mass conservation $3W_0 + W_2 + W_8 = 3$ (Lem.~\ref{lem:mass});
feasibility $\Rightarrow$ density bound (Krasikov--Lagarias \cite{KL03}). \\
\hline
\textbf{Verified} (finite computation) &
Certificates $k = 13, 15, 17, 19, 20$ with
$\gam = 0.8624 / 0.8805 / 0.8953 / 0.9069 / 0.9146$ (Thm.~\ref{thm:certs});
identity \eqref{eq:minloss} to $8$ decimals at $k = 5\ldots11$;
linear edge spectrum $\{1\} \cup \{|\mu| = 1/4\}$ at $j = 4, 5, 6$
(Rem.~\ref{rem:linear}); exact pointwise transport identity
(residual $2.8 \times 10^{-4}$); geometric-$w$ cylinder densities exact
(Terras \cite{Ter76}). \\
\hline
\textbf{Measured} &
Tempering law $\alpha_k$, $R^2$ (Table~\ref{tab:tempering});
$1 - \alpha_k = \CV_{\mathrm{res}}(k)$; constant $0.698$;
edge-rate ratios $0.824 \to 0.908$ (Cor.~\ref{cor:ratios});
closure \eqref{eq:lattice} and coefficients (Table~\ref{tab:theta});
$a + c = 0.9955$ at $k = 20$; fine-end saturation (Prop.~\ref{prop:cv1});
out-of-sample hits at $k = 19, 20$. \\
\hline
\textbf{Conjectured} &
Conjecture T: $\alpha_k \to 1$;
Open Problem~\ref{open:drift}: $a < c$ uniformly ($\theta_\infty < 1$). \\
\hline
\end{tabular}
\end{center}

\medskip

The picture that emerges is simple to state. The Krasikov--Lagarias certificate
--- the best known rigorous handle on the density of $3x+1$ convergence --- is
not an unstructured numerical artifact: it is, to $R^2 > 0.99$ at a billion
classes, one exactly solvable measure raised to one scalar power. The distance
of that power from $1$ \emph{is} the distance of the density exponent from $1$,
through an exchange rate that is an explicit derivative, $1/\log(4/3)$. The
defining constant of the problem pins the certificate's heterogeneity cascade
to the conservative line; the source of heterogeneity saturates; damping is
spectrally overwhelming at the linear level. What is not yet proved is one
inequality about the direction of drift along the conservative line.

The framework is falsifiable at the next depth: at $k = 21$ it predicts
$\gam \approx 0.918$, $\theta \approx 0.850$, and a tempering fit with
$R^2 > 0.9977$.

\subsection*{Reproducibility}
All certificates, verification logs, and the scripts generating every number
in this paper are archived in the research dossier accompanying this work
(directories \texttt{certificates/}, \texttt{scripts/}, \texttt{findings/}).

\subsection*{Acknowledgment of method}
The computations, fits, and drafts were produced in an extended
human--AI collaboration; every proved statement above has a self-contained
proof in the text, independent of that provenance.

\begin{thebibliography}{9}

\bibitem{KL03}
I.~Krasikov and J.~C.~Lagarias,
\emph{Bounds for the 3x+1 problem using difference inequalities},
Acta Arith. \textbf{109} (2003), no.~3, 237--258.
arXiv:math/0205002.

\bibitem{Lag85}
J.~C.~Lagarias,
\emph{The 3x+1 problem and its generalizations},
Amer. Math. Monthly \textbf{92} (1985), 3--23.

\bibitem{Tao19}
T.~Tao,
\emph{Almost all orbits of the Collatz map attain almost bounded values},
Forum Math. Pi \textbf{10} (2022), e12.
arXiv:1909.03562 (2019).

\bibitem{Ter76}
R.~Terras,
\emph{A stopping time problem on the positive integers},
Acta Arith. \textbf{30} (1976), 241--252.

\end{thebibliography}

\end{document}
